%% ---------------------------------------------------------------------
%%  I Radicali: Schema Visivo — la scheda del modello
%%
%%  Riproduce la scheda mandata dall'operatore il 25/11/2026: due sezioni
%%  (campo di esistenza, operazioni), le forme al posto dei numeri.
%%  Qui non c'è grafica: tutto sta in ../schede.sty.
%% ---------------------------------------------------------------------
\documentclass[10pt]{article}
\usepackage{schede}

%% La scheda intera sta in UNA pagina, e non solo sullo schermo: testata e
%% piede dentro la zona che ogni stampante sa stampare (il piede di norma
%% sta a 5,4 mm dal bordo della carta, e a stampa finiva su un secondo
%% foglio). IL RESPIRO RESTA QUELLO DI NORMA: misurato il 26/11 su questa
%% scheda, il corpo finisce a 793 px di 1170 (100 dpi) e sopra il piede
%% restano 78 mm di bianco — \schedaCompatta toglieva 10 mm di spazi a
%% fronte di 78 liberi, e schiacciava le formule senza che servisse. La
%% pagina è una con i margini sicuri da soli (pdfinfo: Pages: 1).
\schedaMarginiSicuri

\schedaTitolo{I Radicali: Schema Visivo}{Campo di esistenza \& operazioni base}
\schedaPiede{Schema di Matematica — I Radicali}

\begin{document}

%% =====================================================================
\schedaSezione{1}{Campo di Esistenza (CE)}
%% =====================================================================
\begin{schedaTabella}{Tipo di radice}{Schema visivo}{Esempio}

  \riga{\voce{Indice pari}{(es. $\sqrt{\ }$, $\sqrt[4]{\ }$, $\sqrt[6]{\ }$)}}
       {\cePari{\formaQuadrato}}
       {\cePari[4]{\formaQuadrato}}

  \riga{\voce{Indice dispari}{(es. $\sqrt[3]{\ }$, $\sqrt[5]{\ }$, $\sqrt[7]{\ }$)}}
       {\ceDispari[3]{\formaQuadrato}}
       {\ceDispari[5]{\formaQuadrato}}

  \riga{\voce{Più radici pari}{(sistema di condizioni)}}
       {$\radice{\formaQuadrato}+\radice{\formaTriangolo}\;\freccia\;
         \sistema{\formaQuadrato\ge 0\\ \formaTriangolo\ge 0}$}
       {$\radice{\formaQuadrato}+\radice[4]{\formaTriangolo}\;\freccia\;
         \sistema{\formaQuadrato\ge 0\\ \formaTriangolo\ge 0}$}

\end{schedaTabella}

%% =====================================================================
\schedaSezione{2}{Operazioni con i Radicali (casi semplici)}
%% =====================================================================
\begin{schedaTabella}{Operazione}{Schema visivo (stesso indice $n$)}{Esempio}

  \riga{\voce{Somma e sottrazione}{(solo radicali simili)}}
       {$a\,\radice[n]{\formaQuadrato}\;\pm\;b\,\radice[n]{\formaQuadrato}
         \;\freccia\;(a\pm b)\,\radice[n]{\formaQuadrato}$\par\vspace{3pt}
        {\footnotesize\color{schedeGrigio}Con radici diverse:
         $a\radice{\formaQuadrato}+b\radice{\formaTriangolo}$ resta così}}
       {$3\radice{\formaQuadrato}+5\radice{\formaQuadrato}
         =8\radice{\formaQuadrato}$\par\vspace{3pt}
        $4\radice{\formaQuadrato}+2\radice{\formaTriangolo}
         -\radice{\formaQuadrato}$\par
        $=3\radice{\formaQuadrato}+2\radice{\formaTriangolo}$}

  \riga{\voce{Prodotto}{(moltiplicazione)}}
       {$\radice[n]{\formaQuadrato}\cdot\radice[n]{\formaTriangolo}
         \;\freccia\;\radice[n]{\formaQuadrato\cdot\formaTriangolo}$}
       {$\radice[3]{\formaQuadrato}\cdot\radice[3]{\formaTriangolo}
         =\radice[3]{\formaQuadrato\cdot\formaTriangolo}$}

  \riga{\voce{Divisione}{(quoziente)}}
       {$\radice[n]{\formaQuadrato}:\radice[n]{\formaTriangolo}
         \;\freccia\;\radice[n]{\formaQuadrato:\formaTriangolo}$\par\vspace{3pt}
        {\footnotesize\color{schedeGrigio}con $\formaTriangolo\neq 0$}}
       {$\radice[4]{\formaQuadrato}:\radice[4]{\formaTriangolo}
         =\radice[4]{\formaQuadrato:\formaTriangolo}$}

  \riga{\voce{Potenza di radice}{}}
       {$\left(\radice[n]{\formaQuadrato}\right)^{k}
         \;\freccia\;\radice[n]{\formaQuadrato^{k}}$}
       {$\left(\radice[3]{\formaQuadrato}\right)^{2}
         =\radice[3]{\formaQuadrato^{2}}$\par\vspace{3pt}
        $\left(2\radice[3]{\formaQuadrato}\right)^{2}
         =4\radice[3]{\formaQuadrato^{2}}$}

  \riga{\voce{Radice di radice}{}}
       {$\radice[m]{\radice[n]{\formaQuadrato}}
         \;\freccia\;\radice[m\cdot n]{\formaQuadrato}$}
       {$\radice[3]{\radice[2]{\formaQuadrato}}
         =\radice[3\cdot 2]{\formaQuadrato}
         =\radice[6]{\formaQuadrato}$}

\end{schedaTabella}

\end{document}
